\section{Results}
\label{sec:results}

\subsection{Overall Performance}
\label{sec:overall-pass}

\begin{table*}[t]
\centering
\caption{Overall pass rates. \qwenshort{} covers 626 valid records (426 L0 + 200 ablation) after excluding 82 records involving 9 \knownfail{} specs. \gptnew{} evaluation is in progress.}
\label{tab:overall-pass}
\small
\begin{tabular}{@{}lrrrrrrl@{}}
\toprule
Model & \pass{} & \buildfail{} & \runfail{} & \verifyfail{} & EXTR\_FAIL & Total & Rate [95\% Wilson CI] \\
\midrule
\qwen{} & 230 & 245 & 121 & 29 & 1 & 626 & 36.7\% [33.1\%, 40.6\%] \\
\gptnew{} & \tbd{} & \tbd{} & \tbd{} & \tbd{} & \tbd{} & \tbd{} & \tbd{} \\
\bottomrule
\end{tabular}
\end{table*}

Kernel-centric evaluation reveals translation capability that repository-level evaluation obscures. \qwenshort{} passes 230 of 626 valid records, or 36.7\%~[33.1\%, 40.6\%], across five suites, all augmentation levels, and eight directions including case studies (Table~\ref{tab:overall-pass}).

The comparison to repository-level work is the main methodological point. ParEval-Repo reports 0\% pass@1 for all models on applications larger than 133~SLoC~\cite{ParEvalRepo2025}. \parbench{} evaluates many kernels above that size but fixes the build and integration scaffold, yielding 36.7\% overall and 40.3\%~[29.7\%, 51.8\%] on \qwenshort{} CUDA-to-OpenMP at L0. This does not make the task easy; it separates kernel translation failures from full-repository reconstruction failures.

\subsection{pass@$k$ Analysis}
\label{sec:passk-analysis}

\begin{table}[t]
\centering
\caption{pass@$k$ results for \qwenshort{} (L0 only, temperature~0.7, three independent samples per task, 142 unique source-target pairs).}
\label{tab:pass-at-k-main}
\small
\begin{tabular}{@{}lr@{}}
\toprule
Metric & Value [95\% CI] \\
\midrule
pass@1 & 23.9\% [17.9\%, 30.0\%] \\
pass@3 & 35.2\% [27.3\%, 43.1\%] \\
\midrule
Always pass (3/3) & 19 tasks (13.4\%) \\
Noisy (1--2/3) & 31 tasks (21.8\%) \\
Hard fail (0/3) & 92 tasks (64.8\%) \\
\bottomrule
\end{tabular}
\end{table}

Stochastic sampling provides limited rescue. Of 142 unique L0 source-target pairs, 92 (64.8\%) are hard failures with no passing sample across three independent attempts (Table~\ref{tab:pass-at-k-main}). Only 19 tasks (13.4\%) pass in all three samples, while 31 (21.8\%) are noisy---passing in some samples but not others. The gap between pass@1 (23.9\%) and pass@3 (35.2\%) is modest, indicating that additional samples yield diminishing returns. Most failures are structural at the kernel-direction level rather than stochastic.

\subsection{Direction Dependence}
\label{sec:direction-analysis}

\begin{table*}[t]
\centering
\caption{Pass rates by direction (L0 records, including all three samples per task). OpenMP-target rows are case studies rather than part of the standard six-direction matrix.}
\label{tab:direction-rates}
\small
\begin{tabular}{@{}lrcl@{}}
\toprule
Direction & \qwenshort{} & $n$ & Note \\
\midrule
CUDA$\to$OMP & 40.3\% [29.7\%, 51.8\%] & 72 & Primary \\
OMP$\to$OpenCL & 33.3\% [22.0\%, 47.0\%] & 51 & \\
OMP$\to$CUDA & 25.0\% [16.4\%, 36.1\%] & 72 & \\
OpenCL$\to$OMP & 9.8\% [4.3\%, 21.0\%] & 51 & \\
CUDA$\to$OpenCL & 7.0\% [2.8\%, 16.7\%] & 57 & \\
OpenCL$\to$CUDA & 0.0\% [0.0\%, 6.3\%] & 57 & \\
\midrule
OMP-target$\to$CUDA & 66.7\% [46.7\%, 82.0\%] & 24 & Case study \\
CUDA$\to$OMP-target & 0.0\% [0.0\%, 13.8\%] & 24 & Case study \\
\midrule
\gptnew{} (all directions) & \tbd{} & \tbd{} & \\
\bottomrule
\end{tabular}
\end{table*}

Performance is strongly direction-dependent (Table~\ref{tab:direction-rates}). For \qwenshort{}, CUDA-to-OpenMP is the strongest standard direction at 40.3\%, followed by OMP-to-OpenCL at 33.3\% and OMP-to-CUDA at 25.0\%. OpenCL-to-CUDA yields 0\% across all 57 L0 records. This pattern reflects the structural asymmetry of the APIs: translating from CUDA to OpenMP removes explicit device memory management and thread launch syntax, whereas translating into CUDA or OpenCL requires introducing those details coherently. OpenCL translations additionally face a separate host/kernel source model and address-space qualifiers that surface compilation errors only at runtime through \texttt{clBuildProgram()}.

The OMP-target case study shows a striking asymmetry: OMP-target-to-CUDA passes at 66.7\% while CUDA-to-OMP-target passes at 0\%. OMP-target code uses directive-based offloading that maps straightforwardly to CUDA's explicit device model, but generating correct \texttt{\#pragma omp target} directives with proper data-mapping clauses from CUDA kernels appears beyond current model capability.

At the kernel level, difficulty is heterogeneous. Table~\ref{tab:per-kernel-main} shows the five easiest and five hardest kernels under \qwenshort{}. The easier kernels tend to be regular stencil, graph, or solver kernels with structured loop nests; the hard set includes irregular control flow, domain-specific data structures, brittle synchronization, and multi-file dependencies. The full kernel table is in Appendix Table~\ref{tab:per-kernel-full}.

\begin{table}[t]
\centering
\caption{Representative per-kernel pass rates for \qwenshort{} across all directions and augmentation levels.}
\label{tab:per-kernel-main}
\small
\begin{tabular}{@{}llrrrr@{}}
\toprule
Suite & Kernel & Total & \pass{} & Fail & Rate \\
\midrule
\multicolumn{6}{l}{\textit{Easiest kernels}} \\
HeCBench & stencil1d & 15 & 15 & 0 & 100.0\% \\
HeCBench & floydwarshall & 20 & 16 & 4 & 80.0\% \\
HeCBench & heat2d & 20 & 15 & 5 & 75.0\% \\
HeCBench & iso2dfd & 20 & 15 & 5 & 75.0\% \\
Rodinia & hotspot & 30 & 21 & 9 & 70.0\% \\
\midrule
\multicolumn{6}{l}{\textit{Hardest kernels}} \\
Rodinia & gaussian & 10 & 0 & 10 & 0.0\% \\
Rodinia & heartwall & 30 & 0 & 30 & 0.0\% \\
Rodinia & myocyte & 30 & 0 & 30 & 0.0\% \\
RSBench & rsbench & 30 & 0 & 30 & 0.0\% \\
XSBench & xsbench & 30 & 0 & 30 & 0.0\% \\
\bottomrule
\end{tabular}
\end{table}

\subsection{Failure Taxonomy}
\label{sec:failure-taxonomy}

\begin{figure}[t]
\centering
\includegraphics[width=\columnwidth]{f4_failure_taxonomy_qwen}
\caption{\qwenshort{} outcome taxonomy across 626 valid records. Build-stage failures are the largest non-pass category. \gptnew{} figures are in Appendix~\ref{sec:appendix-gpt-figures}.}
\label{fig:failure-taxonomy}
\end{figure}

Build-stage adaptation is the dominant bottleneck. \buildfail{} accounts for 245 of 626 records (39.1\%) and 61.9\% of all failures (Table~\ref{tab:overall-pass}, Figure~\ref{fig:failure-taxonomy}). These failures indicate incomplete API-surface adaptation: retained CUDA identifiers in OpenMP translations, missing OpenCL address-space qualifiers, or mismatched target-side dependencies.

\verifyfail{} is smaller but important: 29 of 626 records (4.6\%) compile and run but fail correctness. These cases are the strongest evidence that compilation alone is insufficient: exit-code-only or compile-only checks would misclassify them as successful translations. \runfail{} accounts for 121 records (19.3\%), and a single \extractionfail{} (0.2\%) occurred where the LLM response could not be parsed into target files. Detailed failure case studies appear in Appendix~\ref{sec:appendix-c}.

\subsection{Augmentation Robustness}
\label{sec:augmentation-robustness}

\begin{table}[t]
\centering
\caption{Augmentation pass rates on the balanced 12-kernel CUDA-to-OpenMP subset present at all five levels.}
\label{tab:aug-balanced}
\small
\begin{tabular}{@{}lrrl@{}}
\toprule
Level & Pass / Total & Rate & Kernels with any pass \\
\midrule
L0 & 29/36 & 80.6\% & 12/12 (100\%) \\
L1 &  9/12 & 75.0\% &  9/12 ~~(75\%) \\
L2 & 10/12 & 83.3\% & 10/12 ~~(83\%) \\
L3 &  9/12 & 75.0\% &  9/12 ~~(75\%) \\
L4 & 10/12 & 83.3\% & 10/12 ~~(83\%) \\
\bottomrule
\end{tabular}
\end{table}

If baseline success were driven primarily by surface-form memorization, augmented source variants should degrade pass rates. On a balanced 12-kernel CUDA-to-OpenMP subset present at all five augmentation levels, pass rates remain stable between 75\% and 83\% (Table~\ref{tab:aug-balanced}). L0 has higher counts because the canonical evaluation generates three samples per task (36 records for 12 kernels), while L1--L4 each generate one record per kernel.

The sample sizes at L1--L4 (12 records per level) are too small for a formal monotonic trend test. We therefore report descriptive rates only and note the absence of evidence for systematic degradation. The observed stability is consistent with translation driven by structural understanding of the parallel APIs rather than rote reproduction of memorized source patterns, though the small sample cannot rule out moderate effects. Full per-kernel augmentation heatmaps appear in Appendix Figure~\ref{fig:aug-heatmap}.
